%% kapa: exam=1 round=34 year=2023 period=1 form=A subject=civil_law subject_ko=민법 seq=14 number=14
%% source: Q-Net 감정평가사 기출문제 (cst00309) artlSeq=5211789 file=raw/1008/5211789/2023년도 제34회 감정평가사 1차시험 1교시.pdf
%% conversion: kordoc 4.24.51 -> tools/gen1.py (auto); review=proofread-sonnet; flags=none
\begin{question}{14}
법률행위의 대리에 관한 설명으로 옳지 않은 것은? (다툼이 있으면 판례에 따름)
\begin{choices}
  \item 무권대리인의 상대방에 대한 책임은 대리권의 흠결에 관하여 대리인에게 귀책사유가 있는 경우에만 인정된다.
  \item 민법 제124조에서 금지하는 자기계약이 행해졌다면 그 계약은 유동적 무효이다.
  \item 행위능력자인 임의대리인이 성년후견개시 심판을 받아 제한능력자가 되면 그의 대리권은 소멸한다.
  \item 대리인이 수인인 경우, 법률 또는 수권행위에서 다른 정함이 없으면 각자가 본인을 대리한다.
  \item 상대방 없는 단독행위의 무권대리는 특별한 사정이 없는 한 확정적 무효이다.
\end{choices}
\end{question}
\answer{1}
